% Part: many-valued-logic
% Chapter: three-valued-logics
% Section: introduction

\documentclass[../../../include/open-logic-section]{subfiles}

\begin{document}

\olfileid{mvl}{thr}{int}

\olsection{Pambuka}

Yen kita nambahake siji nilai maneh, yaiku~$\Undef$, marang
$\True$ lan $\False$, kita oleh logika telung nilai. Senajan
nilai kabenerane mung tambah siji, cara netepake fungsi kabeneran
kanggo $\lnot$, $\land$, $\lor$, lan $\lif$ akeh banget.
Sawijining logika bisa nggunakake kombinasi apa wae saka fungsi
kabeneran kasebut. Kita uga bisa milih mung $\True$ minangka
nilai pinilih, utawa milih $\True$ lan~$\Undef$ bebarengan.

Ing kene kita ngandharake sawetara logika telung nilai kang
paling kawentar, motivasine, lan sawetara sipate.

\end{document}
